\documentclass[aspectratio=169,11pt]{beamer}
\usepackage[utf8]{inputenc}
\usepackage[T1]{fontenc}
\usepackage[brazilian]{babel}
\usepackage{lmodern}
\usepackage{tikz}
\usetikzlibrary{arrows.meta,positioning}
\usepackage{booktabs}

% Fonte: fase1_proposta.tex. Roteiro principal: aproximadamente 4min50s.
% Nome para submissão:
% fase1_slides_rafael-attilio-agricola_jociclelio-castro-macedo-junior.pdf
\newcommand{\segundoautor}{Jociclelio Castro Macedo Junior}
\newcommand{\segundora}{231722}

\definecolor{navy}{HTML}{17324D}
\definecolor{teal}{HTML}{007F82}
\definecolor{paper}{HTML}{F1F6F8}
\definecolor{orange}{HTML}{B75A18}
\setbeamercolor{normal text}{fg=navy,bg=white}
\setbeamercolor{structure}{fg=teal}
\setbeamercolor{frametitle}{fg=navy}
\setbeamercolor{title}{fg=navy}
\setbeamercolor{subtitle}{fg=teal}
\setbeamercolor{block title}{fg=white,bg=teal}
\setbeamercolor{block body}{fg=navy,bg=paper}
\setbeamercolor{alerted text}{fg=orange}
\setbeamerfont{title}{size=\LARGE,series=\bfseries}
\setbeamerfont{frametitle}{size=\Large,series=\bfseries}
\setbeamerfont{block title}{size=\normalsize,series=\bfseries}
\setbeamertemplate{navigation symbols}{}
\setbeamertemplate{itemize items}[circle]
\setbeamersize{text margin left=9mm,text margin right=9mm}
\setbeamertemplate{footline}{%
  \hspace{9mm}\textcolor{teal}{\scriptsize EduAgent-OS}\hfill
  \textcolor{navy}{\scriptsize MO810A/MC959A -- Fase 1\quad|\quad
  \insertframenumber/\inserttotalframenumber}\hspace{9mm}\vspace{3mm}}
\setbeamertemplate{frametitle}{\vspace{3mm}\insertframetitle\par
  \vspace{1mm}{\color{teal}\hrule height 0.6pt}\vspace{2mm}}
\hypersetup{pdftitle={EduAgent-OS: proposta da fase 1},
  pdfsubject={Tutoria ativa, memória longitudinal e planejamento adaptativo}}

\title{EduAgent-OS}
\subtitle{Sistema Multiagente Local para Tutoria Ativa\\
  e Gestão Longitudinal de Estudos}
\author{Rafael Attilio Agricola -- 249245\\\segundoautor{} -- \segundora}
\institute{Instituto de Computação -- UNICAMP\\
  MO810A/MC959A -- Sistemas Multiagentes com Modelos de Linguagem\\
  Prof. Dr. Julio Cesar dos Reis}
\date{Fase 1 -- Proposta\enspace|\enspace 2026}

\begin{document}

% 15 s
\begin{frame}
  \titlepage
\end{frame}

% 30 s
\begin{frame}{Do material da disciplina à rotina de estudo}
  \begin{columns}[T,onlytextwidth]
    \begin{column}{0.48\textwidth}
      \begin{block}{Cenário: estudante de graduação}
        \begin{itemize}
          \item Múltiplas disciplinas, PDFs e notas de aula.
          \item Avaliações, compromissos e tempo limitado.
          \item Dificuldades que mudam a cada sessão.
        \end{itemize}
      \end{block}
    \end{column}
    \begin{column}{0.48\textwidth}
      \begin{block}{Motivação: ODS 4}
        Apoiar a aprendizagem autônoma conectando o planejamento
        ao que o estudante demonstra compreender.
      \end{block}
      \small Execução local: menor dependência de inferência paga e
      continuidade com materiais disponíveis offline.
    \end{column}
  \end{columns}
  \vspace{4mm}
  \textbf{Necessidade central:} transformar evidências de aprendizagem
  em decisões sobre \alert{o que estudar e quando revisar}.
  \par\vspace{2mm}
  {\footnotesize ODS 9 e 10: relação complementar; o alcance depende do hardware necessário.}
\end{frame}

% 35 s
\begin{frame}{Problema e questão de pesquisa}
  \textbf{Uma resposta correta a uma dúvida isolada não garante continuidade.}
  \vspace{3mm}
  \begin{enumerate}
    \item \textbf{Tutoria e rotina desconectadas:} dificuldades precisam alterar o plano.
    \item \textbf{Histórico pouco fiel:} distinguir evidência, estimativa antiga e estado atual.
    \item \textbf{Contexto limitado:} documentos e conversas elevam o custo de inferência.
  \end{enumerate}
  \vspace{3mm}
  \begin{block}{Questão de pesquisa}
    Em que medida uma arquitetura multiagente com memória longitudinal
    e gestão seletiva de contexto integra tutoria ativa e replanejamento,
    mantendo qualidade pedagógica e viabilidade temporal com recursos locais limitados?
  \end{block}
\end{frame}

% 30 s
\begin{frame}{Objetivo e recorte do protótipo}
  \begin{block}{Objetivo geral}
    Integrar tutoria ativa, memória longitudinal e planejamento adaptativo:
    as evidências de cada sessão orientam as próximas ações.
  \end{block}
  \begin{columns}[T,onlytextwidth]
    \begin{column}{0.49\textwidth}
      \textbf{Núcleo previsto}
      \begin{itemize}
        \item Organizar tópicos, fontes e prazos.
        \item Tutoria textual, exercícios e flashcards.
        \item Histórico por tópico entre sessões.
        \item Plano semanal com restrições reais.
      \end{itemize}
    \end{column}
    \begin{column}{0.47\textwidth}
      \textbf{Delimitação}
      \begin{itemize}
        \item Inferência, memória e materiais locais.
        \item Web e agenda externa exigem conexão.
        \item Sem notas oficiais, geração de vídeos ou treinamento de modelo fundacional.
      \end{itemize}
    \end{column}
  \end{columns}
\end{frame}

% 45 s
\begin{frame}{Cinco agentes, um estado compartilhado}
  \centering
  \begin{tikzpicture}[
    agent/.style={draw=teal,fill=paper,rounded corners=3pt,
      text width=3.6cm,minimum height=1.12cm,align=center,font=\small},
    hub/.style={agent,fill=navy,text=white,draw=navy,text width=4.5cm},
    flow/.style={<->,>=Stealth,thick,draw=teal}]
    \node[hub] (graph) at (0,0) {\textbf{Orquestração: LangGraph}\\Estado e encaminhamento de tarefas};
    \node[agent] (curador) at (-4.5,1.75) {\textbf{Curador}\\Tópicos, prazos e fontes via RAG};
    \node[agent] (tutor) at (0,1.75) {\textbf{Tutor Adaptativo}\\Perguntas, pistas e exercícios};
    \node[agent] (monitor) at (4.5,1.75) {\textbf{Monitor}\\Evidências e estimativas de domínio};
    \node[agent] (memoria) at (-3,-1.75) {\textbf{Bibliotecário}\\Histórico, origem e ordem temporal};
    \node[agent] (agenda) at (3,-1.75) {\textbf{Planejador}\\Blocos de estudo e revisões};
    \draw[flow] (graph.west) -- (curador.south);
    \draw[flow] (graph.north) -- (tutor.south);
    \draw[flow] (graph.east) -- (monitor.south);
    \draw[flow] (graph.south west) -- (memoria.north);
    \draw[flow] (graph.south east) -- (agenda.north);
  \end{tikzpicture}
  \par\vspace{2mm}
  \raggedright\footnotesize
  \textbf{Inferência proposta:} Bonsai 2 local, compartilhado pelos agentes.\\
  \textbf{Coordenação:} mensagens com tópico, fontes, evidências, prioridade e ação.\\
  \textbf{Ferramentas:} MCP para acesso a dados; A2A a explorar na delegação entre serviços.
\end{frame}

% 35 s
\begin{frame}{Ciclo de uso: a dificuldade muda o plano}
  \begin{enumerate}
    \item \textbf{Organizar:} estudante fornece materiais, avaliações e disponibilidade.
    \item \textbf{Planejar:} propor blocos compatíveis com os compromissos.
    \item \textbf{Recuperar:} trazer objetivos, dificuldades anteriores e fontes.
    \item \textbf{Praticar:} Tutor solicita explicações; Monitor registra evidências.
    \item \textbf{Atualizar:} Bibliotecário consolida o histórico; Planejador ajusta a agenda.
    \item \textbf{Retomar:} estudante recebe a atualização que orienta a próxima sessão.
  \end{enumerate}
  \vspace{2mm}
  \begin{block}{Exemplo ilustrativo}
    Dificuldade perto de uma prova $\rightarrow$ recomendação de reforço
    $\rightarrow$ busca de horário livre $\rightarrow$ nova atividade na sessão seguinte.
  \end{block}
  \small\textbf{Agenda verificável:} ferramenta determinística checa duração e
  sobreposições. Sem tempo suficiente, o sistema explicita o conflito.
\end{frame}

% 30 s
\begin{frame}{Memória e contexto sob orçamento limitado}
  \begin{columns}[T,onlytextwidth]
    \begin{column}{0.31\textwidth}
      \begin{block}{Dentro da sessão}
        Manter interações recentes e compactar trechos antigos entre chamadas.
        \par\medskip
        \textbf{Originais recuperáveis} para conferir as sínteses.
      \end{block}
    \end{column}
    \begin{column}{0.31\textwidth}
      \begin{block}{Entre sessões}
        Recuperar histórico pertinente à disciplina e ao tópico.
        \par\medskip
        Preservar \textbf{evidência, origem, incerteza e temporalidade}.
      \end{block}
    \end{column}
    \begin{column}{0.31\textwidth}
      \begin{block}{Por tarefa}
        Carregar somente a \textit{skill} necessária: explicar, criar flashcards ou questionários.
        \par\medskip
        Comparar \textbf{instruções completas e reduzidas}.
      \end{block}
    \end{column}
  \end{columns}
  \vspace{4mm}
  \textbf{Hipótese a testar:} recuperação seletiva e compactação podem reduzir
  custos preservando a continuidade e a qualidade do acompanhamento.
\end{frame}

% 30 s
\begin{frame}{Estado da arte: elementos e adaptações}
  \small
  \renewcommand{\arraystretch}{1.05}
  \begin{tabular}{@{}p{2.65cm}p{4.25cm}p{6.25cm}@{}}
    \toprule
    \textbf{Trabalho} & \textbf{Elemento aproveitado} & \textbf{Adaptação ao EduAgent-OS}\\
    \midrule
    MetaGPT (2024) & Papéis e artefatos estruturados & Coordenar tópicos, evidências e restrições.\\
    MultiTutor (2025) & Agentes de suporte educacional & Conectar tutoria à rotina e ao histórico.\\
    MemGPT (2023) & Contexto ativo e memória externa & Recuperar progresso e dificuldades por tópico.\\
    Voyager (2023) & Habilidades reutilizáveis & Biblioteca de procedimentos pedagógicos.\\
    ToRA (2024) & Raciocínio com ferramentas & Verificar restrições de agenda deterministicamente.\\
    \bottomrule
  \end{tabular}
  \vspace{3mm}
  \textbf{Diferencial proposto:} integrar tutoria ativa, memória longitudinal
  e replanejamento orientado por evidências em execução local.
\end{frame}

% 40 s. Referências no slide seguinte: apoio, fora do roteiro oral.
\begin{frame}{Como avaliar a proposta nas próximas fases?}
  \small
  \begin{columns}[T,onlytextwidth]
    \begin{column}{0.48\textwidth}
      \textbf{Comparação controlada}
      \begin{itemize}
        \item Multiagente versus assistente único.
        \item Mesmo modelo, materiais e ferramentas; recursos controlados.
        \item Sequências de sessões e cenários de estudo.
        \item Ablações de memória, compactação e seleção de \textit{skills}.
      \end{itemize}
    \end{column}
    \begin{column}{0.48\textwidth}
      \textbf{Dimensões observadas}
      \begin{itemize}
        \item Qualidade e fundamentação das respostas.
        \item Fidelidade da memória entre sessões.
        \item Restrições de agenda e delegações.
        \item Tokens, latência e memória computacional.
      \end{itemize}
    \end{column}
  \end{columns}
  \vspace{1mm}
  \begin{block}{Resultado esperado}
    Protótipo do ciclo completo e registros para avaliar cada mecanismo.
  \end{block}
  \footnotesize A avaliação inicial mede o sistema; ganhos de aprendizagem humana
  exigem investigação adicional.
\end{frame}

\begin{frame}{Referências e elaboração}
  \footnotesize
  \begin{thebibliography}{9}
    \setlength{\itemsep}{1mm}
    \bibitem{metagpt} Hong et al. (2024). \textit{MetaGPT: Meta Programming for A Multi-Agent Collaborative Framework}. ICLR.
      \href{https://arxiv.org/abs/2308.00352}{arXiv:2308.00352}.
    \bibitem{multitutor} Sun e Tai (2025). \textit{MultiTutor: Collaborative LLM Agents for Multimodal Student Support}. PMLR, 273:174--190.
      \href{https://proceedings.mlr.press/v273/sun25a.html}{Artigo}.
    \bibitem{memgpt} Packer et al. (2023). \textit{MemGPT: Towards LLMs as Operating Systems}.
      \href{https://arxiv.org/abs/2310.08560}{arXiv:2310.08560}.
    \bibitem{voyager} Wang et al. (2023). \textit{Voyager: An Open-Ended Embodied Agent with Large Language Models}.
      \href{https://arxiv.org/abs/2305.16291}{arXiv:2305.16291}.
    \bibitem{tora} Gou et al. (2024). \textit{ToRA: A Tool-Integrated Reasoning Agent for Mathematical Problem Solving}. ICLR.
      \href{https://arxiv.org/abs/2309.17452}{arXiv:2309.17452}.
  \end{thebibliography}
  \vspace{1mm}
  \textbf{Uso de IA generativa:} proposta elaborada com auxílio de Claude Haiku 4.5,
  Muse Spark 1.2 Free e GPT-6-Astra; adaptação para slides com GPT-6-Astra.
  Revisão final e responsabilidade pelo conteúdo cabem aos autores.
\end{frame}

\end{document}
